\chapter{第二十二讲\quad 妨害社会管理犯罪}
\label{lecture:22}
\noindent{\small\hyperref[ref:lecture:22]{查阅本讲扩展规则}}\par

\section{第一节\quad 扰乱公共秩序罪}

\paragraph{法条与分流}
令 \(L\) 表示依法执行职务，\(P\) 表示人民警察，\(A\) 表示暴力袭击，\(H\) 表示以暴力、威胁阻碍。第277条依行为和对象分别检查。第五款的基本分支用 \(L\land P\land A\)，第一款典型分支用依法执行职务的国家机关工作人员与 \(H\)。根据高检发释字〔2025〕1号第1条，直接袭击造成轻微伤以上，以及打砸警车、警械等足以危及人身安全，是识别 \(A\) 的两类路径。以枪支、管制刀具或驾驶机动车撞击等手段严重危及民警人身安全的，另审第五款加重档。冒充国家机关工作人员招摇撞骗查第279条。第280条按公文、证件、印章的归属及伪造、变造、买卖、盗窃、抢夺、毁灭的法定组合核对，第280条之二冒名顶替要核对法定入学、录用或安置资格。

聚众斗殴（第292条）先查首要分子或积极参加者，再查重伤、死亡特别转化规则。第293条的寻衅滋事行为须落入法定类型并满足相应情节。第291条之二高空抛物须查情节严重，同时构成其他犯罪按处罚较重规定。第294条的黑社会性质组织用组织、经济、行为、危害性四组条件。第303条分别看赌博、开设赌场及组织参与境外赌博的法定分支。虚假危险物质、恐怖信息、险情疫情灾情警情分别对照第291条之一的对应款。

\paragraph{推导}
警察乙正在依法抓捕，甲为阻止抓捕当面挥拳击打乙，造成轻微伤。执法凭证和行动事实给出 \(P,L\)，挥拳与伤情依司法解释给出 \(A\)。依第277条第五款，\(P\land L\land A\) 推出袭警罪的基本候选。若将挥拳换成口头反对且无威胁、无暴力，\(A,H\) 都失效，原推导不能继续。

\paragraph{边界与邻近模型}
组织考试作弊罪的考试须是法律规定的国家考试（第284条之一）。学校内部周测不是因「考试」二字就符合。计算机案件先标出行为是侵入、取数据、控制、破坏、建网发布，还是提供帮助。第285条第一款限于法定特殊领域系统。其他系统非法获取数据、非法控制等审查第二款及严重情节。第286条查破坏后果。第287条之一、之二分别核对非法利用信息网络的行为类型与明知他人利用信息网络实施犯罪而提供法定帮助，均须核对严重情节。

\section{第二节\quad 妨害司法罪}

\paragraph{法条与对象坐标}
令 \(O\) 为干扰对象，取值为证言、其他证据、犯罪人、犯罪所得、生效裁判、在押状态。第305条伪证罪还要求刑事诉讼中的证人、鉴定人、记录人或翻译人，对重要情节作虚假证明、鉴定、记录或翻译，意图陷害他人或隐匿罪证。第307条区分阻止作证、指使作伪证与帮助当事人毁灭、伪造证据，后者须情节严重。第307条之一另核对以捏造事实提起民事诉讼，以及妨害司法秩序或者严重侵害他人合法权益。

\paragraph{推导}
盗窃既遂后甲才得知真相，随后替盗窃者藏匿赃物。上游盗窃事实与物品流转记录给出 \(O=\text{犯罪所得}\)。已知赃物性质而藏匿，给出明知与掩隐行为。根据第312条及2025年相关司法解释综合核对入罪，即可在该项事实达标时推出掩饰、隐瞒犯罪所得罪。若只替盗窃者本人提供隐藏住所，且明知对方是犯罪的人，\(O\) 改为犯罪人，应转查第310条窝藏罪。若事前通谋，则进入上游犯罪共犯的审查路径。

\paragraph{边界}
证人在民事诉讼作假，缺失第305条的刑事诉讼条件。第313条要求负有执行义务、有能力执行而拒不执行且情节严重，真实无能力履行不能直接推出拒执罪。第316条的脱逃还要求依法被关押的特定身份。证明链中「人、物、判决、在押」的对象不可混同。

\section{第三节\quad 妨害文物管理罪}

\paragraph{法条与定义}
令 \(W\) 为国家禁止经营的文物，\(M\) 为牟利目的，\(D\) 为倒卖，\(Q\) 为严重情节。第326条的个人分支以 \(W\land M\land D\land Q\) 为核心。第328条另以擅自盗掘具有历史、艺术、科学价值的古文化遗址、古墓葬等法定对象为核心。

\paragraph{推导}
甲以牟利为目的收购国家禁止经营的文物后反复倒卖，已满足严重情节。文物鉴定和经营限制给出 \(W\)，资金与交易记录给出 \(M,D,Q\)，依第326条得倒卖文物罪的候选。若甲在取得该文物之前亲自盗掘有上述价值的古墓葬，应将盗掘行为单列，先核对第328条再处理罪数。

\paragraph{边界}
合法经营的文物不满足 \(W\)。盗掘古墓葬是否成立不由最后是否挖到值钱物品单独决定，须分别确定对象、行为以及完成形态。

\section{第四节\quad 妨害国（边）境管理罪}

\paragraph{法条与模型}
令 \(I\) 为他人意欲非法出入境，\(G\) 为策划、招募、安排等组织行为，\(T\) 为运送。第318条核对 \(I,G\)，第321条核对 \(I,T\) 及运送者相应的认识。第322条是偷越者本人并符合其入罪条件的分支。组织者、运送者与偷越者的要件不能互换。

\paragraph{推导}
甲为偷越者招募人员、定线路、统一安排越境时点，乙仅按指令驾车送人且明知是偷越。甲的组合事实证明 \(G\)，依第318条审查组织分支。乙的知情运送给出 \(T\)，依第321条审查运送分支，再根据乙是否参与组织共同故意核对共犯。仅凭两人都在车上不能将二者同罪处理。

\paragraph{边界}
司机正常接单，证据明确其不知乘客准备偷越，不能仅由客观送达地点推得故意。人数、后果、危险手段以及遇到妨害公务等情形，另审各罪加重和并罚规定。

\section{第五节\quad 危害公共卫生罪}

\paragraph{法条与模型}
第335条的医疗事故罪关注医务人员身份 \(I\)、严重不负责任 \(N\)、就诊人死亡或者严重损害身体健康 \(H\) 以及因果归责 \(C\)。第336条非法行医罪先核对未取得医生执业资格的人非法行医和情节严重，严重损害健康或造成就诊人死亡是另外的结果档。两条的主体资格、基准入罪条件不同。

\paragraph{推导}
有执业资格的医务人员甲值班，明知须在给药前核对药物，却严重违背规范。鉴定及其他证据确定错药造成患者死亡，规范用于预防该种危险。由资格及值班得 \(I\)，由严重违规得 \(N\)，由死因得 \(H,C\)，依第335条得医疗事故罪的候选。

\paragraph{边界}
如果同一病人最终死亡，但死因与医务人员违章行为无关，则 \(C\) 不成立，不能据死亡套用第335条。未取得医生执业资格而非法行医者，即使暂未查明死亡，也应另查第336条的情节严重分支。

\section{第六节\quad 环境犯罪}

\paragraph{法条与模型}
第345条区分盗伐与滥伐。先查林木权属和非法占有目的，再查采伐许可、实际采伐方式和数量。第341条分辨珍贵、濒危野生动物及制品分支、非法狩猎分支、以食用为目的的陆生野生动物相关分支。第338条设 \(N\) 为违反国家规定，\(D\) 为法定排放、倾倒、处置行为，\(Q\) 为严重污染环境。环境污染刑事司法解释的两次修正分别于2026年3月30日、8月15日施行。本节按第二次修正后的文本核对 \(Q\) 的各条路径。

\paragraph{推导}
甲违反国家规定，明知所倾倒的是危险废物仍擅自倾倒，材料性质经鉴定确认。调查确认倾倒三吨以上，所涉时间适用该司法解释。证据分别给出 \(N,D\)，依环境污染解释第1条第（二）项得 \(Q\)。于是客观要件进入如下路径，罪过及责任条件另行核对：
\[
N\land D\land Q\ \xrightarrow{\text{第338条}}\ \text{污染环境罪候选}.
\]
是否积极修复生态，另依该解释第6条审查从宽处理。

\paragraph{边界}
将三吨材料换成经鉴定不属于危险废物的普通材料，该项门槛失效，但仍可调查其他入口。不能用「没有人伤亡」否定所有严重污染情形，也不能用「后来修复」自动抹去已成立的要件。

\section{第七节\quad 毒品犯罪}

\paragraph{法条与模型}
令 \(X\) 为法定毒品，\(K\) 为对毒品性质的认识，\(A\) 取走私、贩卖、运输、制造。第347条规定上述行为不以毒品数量作为统一入罪门槛。第348条非法持有毒品须另查法定数量。第349条区分包庇特定毒品犯罪分子与为犯罪分子窝藏、转移、隐瞒毒品或毒赃，事先通谋时转向上游共犯。第353条查引诱、教唆、欺骗或强迫他人吸毒，第354条查容留，第351条查非法种植及法定入罪分支。

\paragraph{推导}
甲明知粉末是毒品，在国内出售并交给购买者，已经成交，数量低于第348条非法持有的门槛。由鉴定和交易事实得到 \(X,K,A=\text{贩卖}\)，按第347条评价，不能把第348条的持有门槛加入第347条。若只证明甲明知而非法持有，且数量未达第348条标准，无法推出非法持有毒品罪。若贩卖目的尚无证据，也不能从持有直接推得贩卖。

\paragraph{边界}
购买自用、代购、受托携带、专业运输有不同的意思联络和行为事实。移动地点本身不足以把任何持有都认定为第347条运输。毒品数量按第357条核对查证属实的数量，不以纯度折算。第356条的再犯从重另有前罪范围，应当与一般累犯分开。

\section{第八节\quad 卖淫类犯罪}

\paragraph{法条及司法解释}
令 \(G\) 为以招募、雇佣、纠集等手段管理、控制他人卖淫，\(n\) 为受组织人数。组织卖淫的识别用 \(G\land(n\ge3)\)，三人来自2017年相关司法解释第1条。第358条区分组织、强迫与明知他人组织而协助。第359条区分引诱、容留、介绍及引诱幼女卖淫。第360条另查明知患有法定严重性病仍卖淫、嫖娼。

\paragraph{推导}
甲招募并统一排班、控制三名成年卖淫者，乙明知该组织，专职为甲招募、运送人员。招募与控制给出 \(G\)，三人给出人数条件，甲满足组织分支。乙的角色与故意给出法定协助行为，依第358条另审协助组织卖淫罪。换成正常受雇清洁且证据表明不知情，不得保留乙的协助故意。

\paragraph{边界}
偶发介绍不等于持续管理控制。行为人对被组织卖淫者另作引诱、容留、介绍，应按2017年解释第3条核对择重规定。只因同在一处工作不能推出组织者与受雇人的共同罪名。

\section{第九节\quad 关于淫秽物品的犯罪}

\paragraph{法条与模型}
\(X\) 表示第367条意义上的淫秽物品，\(A\) 表示法定制作、复制、出版、贩卖、传播行为，\(P\) 表示牟利目的。第363条第一款的分支核对 \(X,A,P\) 和司法解释适用标准。第364条第一款的传播分支核对 \(X\)、传播与情节严重，不要求牟利目的。组织播放的规定另审，不能用「制作」替换「传播」而保持结论。

\paragraph{推导}
甲为收费获利，经营付费网站传播经认定的淫秽物品，且达到适用入罪标准。从收费经营事实得 \(P\)，内容和传播记录得 \(X,A\)，按第363条第一款审查。若仅改成无牟利目的的传播，\(P\) 失效，但若传播达到严重情节，仍可按第364条第一款核对。

\paragraph{边界}
关于人体生理、医学知识的科学著作和有艺术价值的作品应依第367条审查排除条件。不同违禁品的「贩卖」「出售」「买卖」用语必须逐条看对象和行为。不能先给这些词一个跨罪名的统一含义，再凭字面判处数罪。

\paragraph{闭卷记忆}
\begin{enumerate}
\item 画九格，依次标记「公共、司法、文物、国境、卫生、环境、毒品、卖淫、淫秽」。把身份、对象、行为和严重程度填在各格。
\item 把司法格拆成「人、物、证、判、押」。分别对应犯罪人、赃款赃物、证据、裁判和在押状态，标出参与发生在上游之前还是之后。
\item 对卫生记「资格、严重违规、后果、归责」，对毒品记「对象、认识、动作、数量」。先判断分支，再问该分支把数量放在入罪还是量刑。
\end{enumerate}
口诀「\textbf{先抓对象，再看动作。资格、时间、数量各归其格}」。

\paragraph{闭卷自测与答案}
\begin{enumerate}
\item 盗窃已经完成，乙此前无通谋，事后知情并专门藏匿赃物。第310条和第312条先核对谁？答：物为犯罪所得，应先查第312条的故意、行为和入罪条件。若藏的是犯罪人，改查第310条。
\item 同样是少量毒品，为什么贩卖和持有的推导不同？答：第347条不以数量为统一入罪门槛，第348条设有数量条件。前者仍须查明毒品性质、认识与贩卖行为。
\item 违规排污后积极修复，是否自动不构成污染环境罪？答：先审第338条的构成要件。积极修复可作为从宽因素，免予刑事处罚或者不作为犯罪处理须符合各自条件。
\end{enumerate}
